
\subsection{Pricing basis}
\label{sec:cost-basis}

We price each backend at an August 2026 on-demand rate, and the three
rates do not come from the same place. The TPU figure, \$1.20 per
v5e chip-hour, so \$9.60/hour for an 8-chip pod, is Google Cloud's own
published price \citep{googletpupricing}. The GPU figure,
approximately \$0.35/hour for an NVIDIA T4 instance, is taken from a
third-party survey of Google Cloud prices carried out in the same
month. The CPU figure, approximately \$0.19/hour for an instance
comparable to \texttt{n2-standard-4}, is our own estimate, for which we
recorded no dated source. It is the weakest input here, and we use it
only for the CPU row. The catalogue marks all three as input
assumptions, not measurements. These are list prices
for a hypothetical always-on rental, not the cost we were actually
billed: our TPU cluster was a Stanford course allocation, and our
Colab CPU/GPU sessions ran on shared, non-dedicated infrastructure.
We could not obtain a smaller TPU slice than the full 8-chip pod on
our GKE setup, so the single-call TPU cost below bills for the whole
pod even though only one chip does work. The chip-count analysis in
Section~\ref{sec:cost-scaling} instead prices slices of 1, 2, 4 and 8
chips at the same per-chip rate, which prices \emph{hypothetical}
slices: 2- and 4-chip subsets failed to initialize on our setup with a
libtpu topology error, so no smaller slice was ever billed.
The \$1.20 rate was still listed unchanged when we re-checked Google's
pricing page on 2026-09-17.

\subsection{Cost per prediction, single call}
\label{sec:cost-baseline}

Combining these rates with the steady-state throughput measured in
Section~\ref{sec:results-hardware} gives the per-backend cost of 1{,}000
predictions in Table~\ref{tab:cost-baseline}.

\begin{table}[H]
\centering\small
\caption{Cost per 1{,}000 single-call predictions at on-demand list
prices, from August steady-state throughput. The TPU row bills the full
8-chip pod.}
\label{tab:cost-baseline}
\begin{tabular}{lrr}
\toprule
Backend & Throughput (pred/h) & \$ / 1{,}000 predictions \\
\midrule
CPU                          & 17.0   & \$11.19 \\
GPU (T4)                     & 275    & \$1.27  \\
TPU pod, 1 of 8 chips active & 7{,}660 & \$1.25  \\
\bottomrule
\end{tabular}
\end{table}

The TPU pod and the GPU land at almost the same \$/1k despite the pod
costing 27.4$\times$ more per hour: the chip is far faster than the
GPU, 7{,}660 against 275 predictions per hour, but on this single-chip
workload the pod bills for seven idle chips (87.5\% of the pod, the
same idle-chip finding as Section~\ref{sec:parallelism}), and the two
effects very nearly cancel. If it were possible to rent
one TPU chip alone at the same per-chip rate, the hypothetical cost
would be \$0.157/1k, about 8.1$\times$ cheaper than the GPU. That
makes the pod-billing effect, not the chip's own economics, the reason
the \$1.25 figure looks unremarkable next to the GPU. The near-equality is specific to the August GPU baseline: on the
September rerun the GPU would sit near \$0.64/1k and the pod would be
about twice as expensive (Section~\ref{sec:rh-variability}). The
idle-chip mechanism does not depend on which baseline is used.


Renting the whole pod stops being wasteful once all 8 chips do useful
work. Using the multi-query \texttt{pmap} throughput from
Section~\ref{sec:parallelism} (6.92$\times$ over one chip, 14.718
proteins/second across 8 independent proteins) at the same \$9.60/hour pod
rate gives
$
\$9.60 \,/\, (14.718 \times 3600) \times 1000 \approx \$0.181\text{/1k predictions},
$
within about 15\% of the single-chip hypothetical rate once the pod is
kept busy. The two rates are not strictly commensurable: the
multi-query throughput comes from the rotated-alphabet input family
timed around \texttt{apply}, while the single-chip baseline folds the
fixed synthetic sequence and times \texttt{predict}, which also
includes AlphaFold2's host-side confidence metrics
(Section~\ref{sec:meth-workload}). The like-for-like chip-count
comparison is the 100-residue grid of Section~\ref{sec:cost-scaling}.

\subsection{Cost as a function of chip count}
\label{sec:cost-scaling}

Section~\ref{sec:scaling} finds throughput scaling as
$\mathrm{chips}^{0.963}$, just below linear, which by itself would
predict that cost per prediction is nearly independent of chip count.
We test this directly on the measured chip-count $\times$ sequence-length
grid, not on the fitted formula, in Table~\ref{tab:cost-scaling}. Each
point is the multi-query \texttt{pmap} workload, and each cost charges
only the chips in use (chips $\times$ \$1.20 per hour divided by
measured throughput). This prices a \emph{hypothetical} slice of that
size at the same per-chip rate, unlike Table~\ref{tab:cost-baseline},
which bills the whole 8-chip pod; as Section~\ref{sec:cost-basis}
notes, slices smaller than eight chips could not be obtained on our
setup.

\begin{table}[H]
\centering\small
\caption{Cost per 1{,}000 predictions by chip count and sequence length,
computed from measured throughput on the scaling grid. The last column
is the increase from 1 to 8 chips.}
\label{tab:cost-scaling}
\begin{tabular}{lrrrrr}
\toprule
Length (residues) & 1 chip & 2 chips & 4 chips & 8 chips & 1$\to$8 \\
\midrule
100  & \$0.114 & \$0.121 & \$0.128 & \$0.139 & +23\% \\
250  & \$0.386 & \$0.394 & \$0.399 & \$0.413 & +7\% \\
500  & \$1.004 & \$1.012 & \$1.019 & \$1.036 & +3\% \\
1000 & \$5.051 & \$5.089 & \$5.089 & \$5.109 & +1\% \\
\bottomrule
\end{tabular}
\end{table}

Adding chips is not free: going from 1 to 8 chips raises cost per
prediction by 22.6\% at 100 residues (23\% in the table's integer
rounding, computed from the unrounded throughputs). The penalty shrinks quickly with
sequence length, to 7\% at 250, 3\% at 500 and about 1\% at 1000
residues; the last figure is within the rounding of the one-chip
throughput at that length, which the grid records to only two
significant figures. On this grid, with one run per point, the
near-invariance implied by the fitted exponent therefore holds only for
long sequences. For short sequences, close to
the 118-residue input of our single-query experiments, using more chips
has a real, double-digit cost.

\subsection{Practical implication}
\label{sec:cost-implication}

Chip count should therefore be chosen for latency and throughput
requirements, not to minimize unit cost: the cost penalty for using
more chips is modest, but it is not zero, and it is largest exactly in
the short-sequence regime this study otherwise focuses on. In our
measurements, the dominant cost lever for this workload is not how
many chips are rented but whether they are kept busy. Renting an
8-chip pod
to run a single-chip job wastes 87.5\% of what is being paid for; the
same pod running eight independent proteins in parallel does not cost
more per hour and delivers 6.92$\times$ the predictions, an 86.5\%
parallel efficiency.

All figures in this section use a single steady-state measurement per
backend (Section~\ref{sec:limitations}). We did not recompute the CPU
figures against the September sessions: no session has a better claim
than another to being the reference
(Section~\ref{sec:rh-variability}).
